\documentclass[12pt,a4paper,oneside]{article}
\usepackage[brazil]{babel}
\usepackage[utf8]{inputenc}
\usepackage{olymp}
\usepackage{color}
\usepackage[dvipsnames]{xcolor}
\usepackage{listings}

\setcounter{problem}{0}

\begin{document}

\begin{problem}{Contar}{contar}{2}
 
Dada uma lista de $N$ números inteiros, seu programa deve determinar quantos destes números são maiores ou iguais a um determinado valor $X$.

%\Notes
%
%Observações, se tiver.

\InputFile

A entrada contém três linhas: a primeira contém um valor inteiro $N$, indicando a quantidade de números na lista; a segunda contém os $N$ números da lista, que são números inteiros $V_i$ separados por espaço em branco; e a terceira linha contém um valor inteiro $X$.

Restrições: $1 \leq N \leq 1000$, $-10^6 \leq V_i \leq 10^6$, $-10^6 \leq X \leq 10^6$.



\OutputFile

Seu programa deve gerar apenas uma linha de saída, informando a quantidade de números da lista que são maiores ou iguais a $X$.

%\Notes
%Nesta questão, seu programa deve usar a função \texttt{fatorial} dada %acima exatamente como está, não a modifique! Sua
%tarefa é apenas criar o programa com a função \texttt{main}.

\Example

\begin{example}
\exmp{
10
1 2 3 4 5 6 7 8 9 10
5
}{
6
}%
\end{example}

\begin{example}
\exmp{
10
1 2 3 4 5 6 7 8 9 10
20
}{
0
}%
\end{example}

\begin{example}
\exmp{
1
2
3
}{
0
}%
\end{example}

\begin{example}
\exmp{
4
1 4 6 4
4
}{
3
}%
\end{example}

%Comentários sobre o exemplo, se necessário.

\end{problem}

\end{document}
